\setcounter{section}{3}

\Needspace{5\baselineskip}
\hypertarget{f1bux4e0edualpipe}{%
\section{1F1B与DualPipe}\label{f1bux4e0edualpipe}}

\Needspace{5\baselineskip}
\hypertarget{ux4e00ux5206ux5e03ux5f0fux8badux7ec3ux6846ux67b6ux4e2dux7684ux57faux672cux8badux7ec3ux6b65ux9aa4}{%
\subsection{一、分布式训练框架中的基本训练步骤}\label{ux4e00ux5206ux5e03ux5f0fux8badux7ec3ux6846ux67b6ux4e2dux7684ux57faux672cux8badux7ec3ux6b65ux9aa4}}

\Needspace{4\baselineskip}
\begin{enumerate}
\def\labelenumi{\arabic{enumi}.}
\tightlist
\item
  F（Forward）：前向计算
\end{enumerate}

动作：输入数据 \(x\) 经过当前层的算子（如矩阵乘法 \(Wx\)），输出激活值 \(y\)。

\Needspace{5\baselineskip}
公式：

\[
y = f(x, \theta)
\]

关键任务：计算预测结果，并\textbf{缓存（Checkpointing）}中间层的激活值，留给后续反向传播使用。

在流水线并行中：F 的输出需要发送给下一个（Downstream）GPU 节点。

\Needspace{4\baselineskip}
\begin{enumerate}
\def\labelenumi{\arabic{enumi}.}
\setcounter{enumi}{1}
\tightlist
\item
  B（Backward for Input / Activation Gradient）：激活梯度计算
\end{enumerate}

动作：根据上一层传回的梯度 \(\frac{\partial L}{\partial y}\)，计算对\textbf{当前层输入} \(x\) 的梯度。

\Needspace{5\baselineskip}
公式：

\[
\nabla x = \frac{\partial L}{\partial x} = \frac{\partial L}{\partial y} \cdot \frac{\partial y}{\partial x}
\]

关键任务：它是反向传播的``接力棒''。计算出的 \(\nabla x\) 必须立刻传回给上一个（Upstream）GPU 节点，以便那一层继续计算。

地位：在流水线并行中，B 位于\textbf{关键路径（Critical Path）}上。如果 B 算得慢或传输延迟高，整个流水线都会卡住（产生 Bubble）。

\Needspace{4\baselineskip}
\begin{enumerate}
\def\labelenumi{\arabic{enumi}.}
\setcounter{enumi}{2}
\tightlist
\item
  W（Backward for Weights / Parameter Gradient）：权重梯度计算
\end{enumerate}

动作：根据梯度 \(\frac{\partial L}{\partial y}\)，计算对\textbf{当前层参数} \(\theta\) 的梯度。

\Needspace{5\baselineskip}
公式：

\[
\nabla \theta = \frac{\partial L}{\partial \theta} = \frac{\partial L}{\partial y} \cdot \frac{\partial y}{\partial \theta}
\]

关键任务：计算出的 \(\nabla \theta\) 会被存放在梯度缓冲区中，最终由优化器（Optimizer）用于更新模型权重。

地位：\(W\) \textbf{不在关键路径上}。因为 \(\nabla \theta\) 只需要在这一层本地更新，不需要传给其他 Rank 来解开计算依赖。

\Needspace{5\baselineskip}
\hypertarget{ux4e8c1f1b}{%
\subsection{二、1F1B}\label{ux4e8c1f1b}}

假设我们有 4 个 GPU（Rank 0\textasciitilde3），流水线深度 \(PP = 4\)。

\begin{itemize}
\tightlist
\item
  \textbf{对于 Rank 0 来说}：它处理微批次 \(m_1\) 的前向计算（\(F_1\)）非常快。
\item
  \textbf{但是}：\(F_1\) 算完后，必须传给 Rank 1 \(\to\) Rank 2 \(\to\) Rank 3，在 Rank 3 算出 Loss 后，梯度再从 Rank 3 \(\to\) Rank 2 \(\to\) Rank 1，最后才传回 Rank 0。
\item
  \textbf{结论}：当 Rank 0 拿到 \(m_1\) 的梯度准备算反向（\(B_1\)）时，中间已经过去了一段很长的时间。
\end{itemize}

为了不让 GPU 闲着，Rank 0 在等待 \(B_1\) 的梯度传回来的这段时间里，会继续往后算 \(F_2, F_3, F_4, \dots\)。

\begin{itemize}
\tightlist
\item
  \(F_i\)：代表当前正在处理的第 \(i\) 个微批次的前向。
\item
  \(B_{i-k}\)：代表之前某个\textbf{已经跑完一圈回来}的微批次（第 \(i-k\) 个）的反向。
\end{itemize}

这里的 \(k\) 实际上取决于 GPU 所在的 Rank 位置。

\begin{itemize}
\tightlist
\item
  在流水线末端（如 Rank 3），\(F\) 和 \(B\) 靠得很近，\(k\) 很小。
\item
  在流水线顶端（如 Rank 0），\(F\) 和 \(B\) 隔得很远，\(k\) 很大。
\end{itemize}

\Needspace{5\baselineskip}
``气泡''大小：

\Needspace{5\baselineskip}
假设每个阶段（\(F\) 或 \(B\)）耗时为 1 个单位时间：

\begin{enumerate}
\def\labelenumi{\arabic{enumi}.}
\tightlist
\item
  \(T = 0\)：Rank 0 执行 \(F_1\)。
\item
  \(T = PP - 1\)：经过 \(PP - 1\) 跳，\(F_1\) 到达最后一个 Rank 并完成计算，立即算出 Loss。
\item
  \(T = PP\)：最后一个 Rank 开始执行反向计算 \(B_1\)。
\item
  \(T = PP + (PP - 1)\)：梯度经过 \(PP - 1\) 跳，传回 Rank 0。
\item
  \(T = 2PP - 1\)：Rank 0 终于可以开始执行 \(B_1\)。
\end{enumerate}

结论：在 Rank 0 看来，从 \(F_1\) 结束到 \(B_1\) 开始，中间隔了 \(2PP - 2\) 个时间片。

这里有一个概念上的区分：\textbf{单任务延迟} vs \textbf{流水线总气泡}。

\begin{itemize}
\tightlist
\item
  \textbf{总气泡（Bubble）}：指的是在整个训练任务中，所有 GPU 处于``空闲（Idle）''状态的总时间。
\item
  在 1F1B 中，Rank 0 在等待 \(B_1\) 回来的这段时间内，并没有闲着，而是在疯狂地跑 \(F_2, F_3, \dots, F_{PP}\)。
\item
  真正无法掩盖的空闲时间只有：最开始启动时的``热身''和最后结束时的``收尾''。将所有 GPU 的空闲时间相加再平均，每个 GPU 的平均气泡大小确实约为 \((PP - 1) \times (F + B)\)。
\end{itemize}

\Needspace{5\baselineskip}
\hypertarget{ux4e09dualpipe}{%
\subsection{三、DualPipe}\label{ux4e09dualpipe}}

我们将以 \textbf{4 个 GPU（Rank 0-3）}、流水线并行深度 \(PP = 4\) 为例。

\begin{itemize}
\tightlist
\item
  \textbf{正向流（D1）}：数据从 \(R0 \to R1 \to R2 \to R3\)（逻辑层 L1 \(\to\) L8）。
\item
  \textbf{反向流（D2）}：数据从 \(R3 \to R2 \to R1 \to R0\)（逻辑层 L1 \(\to\) L8，注意是镜像的）。
\item
  \textbf{计算拆分}：每个阶段分为 \(F\)（前向）、\(B\)（算输入梯度，传给邻居）、\(W\)（算权重梯度，本地用）。
\end{itemize}

\Needspace{5\baselineskip}
\hypertarget{ux4e00ux53ccux5411ux70edux8eabux9636ux6bb5}{%
\subsubsection{（一）双向热身阶段}\label{ux4e00ux53ccux5411ux70edux8eabux9636ux6bb5}}

\textbf{\(T1\)：双向同时启动}

\begin{itemize}
\tightlist
\item
  Rank 0：\(F_{D1}(m_1)\)
\item
  Rank 1：-
\item
  Rank 2：-
\item
  Rank 3：\(F_{D2}(n_1)\)
\end{itemize}

\textbf{\(T2\)：数据向中间汇聚}

\begin{itemize}
\tightlist
\item
  Rank 0：\(F_{D1}(m_2)\)
\item
  Rank 1：\(F_{D1}(m_1)\)
\item
  Rank 2：\(F_{D2}(n_1)\)
\item
  Rank 3：\(F_{D2}(n_2)\)
\end{itemize}

\textbf{\(T3\)：中间节点开始处理双向流}

\begin{itemize}
\tightlist
\item
  Rank 0：\(F_{D1}(m_3)\)
\item
  Rank 1：\(F_{D1}(m_2)+F_{D2}(n_1)\)
\item
  Rank 2：\(F_{D2}(n_2)+F_{D1}(m_1)\)
\item
  Rank 3：\(F_{D2}(n_3)\)
\end{itemize}

\textbf{\(T4\)：第一批双向流到达两端}

\begin{itemize}
\tightlist
\item
  Rank 0：\(F_{D1}(m_4)+F_{D2}(n_1)\)
\item
  Rank 1：\(F_{D1}(m_3)+F_{D2}(n_2)\)
\item
  Rank 2：\(F_{D2}(n_3)+F_{D1}(m_2)\)
\item
  Rank 3：\(F_{D2}(n_4)+F_{D1}(m_1)\)
\item
  备注：此时 \(m_1\) 到达 \(R3\)，\(n_1\) 到达 \(R0\)。
\end{itemize}

分析：在 \(T4\) 结束时，\(m_1\) 在 \(R3\) 算完了全部前向（逻辑 L1-L8），\(n_1\) 在 \(R0\) 算完了全部前向。从 \(T5\) 开始，系统将产生第一批梯度并进入复杂的对冲阶段。

注意：m1和n1是完全不同的数据批次。

\Needspace{5\baselineskip}
\hypertarget{ux4e8cux7a33ux6001ux9636ux6bb5}{%
\subsubsection{（二）稳态阶段}\label{ux4e8cux7a33ux6001ux9636ux6bb5}}

\Needspace{5\baselineskip}
在稳态下，一个 GPU 并不是按顺序执行，而是通过两个主要的\textbf{并行块（Candidate Blocks）}来工作。在一个完整的稳态步内，它包含以下两个核心对冲：

\textbf{对冲 A：\(F_{D1}\) 与 \(B_{D2}\) 的完全重叠}

\begin{itemize}
\tightlist
\item
\Needspace{5\baselineskip}
  \textbf{计算内容}：GPU 同时启动两个算子：

  \begin{enumerate}
  \def\labelenumi{\arabic{enumi}.}
  \tightlist
  \item
    正向流新批次的前向计算 \(F_{D1}(m_i)\)。
  \item
    反向流老批次的激活梯度计算 \(B_{D2}(n_j)\)。
  \end{enumerate}
\item
\Needspace{5\baselineskip}
  \textbf{通信掩盖}：

  \begin{itemize}
  \tightlist
  \item
    在计算 \(F_{D1}\) 时，后台接收 \(B_{D2}\) 需要的下游梯度。
  \item
    在计算 \(B_{D2}\) 时，后台发送 \(F_{D1}\) 产生的激活值。
  \end{itemize}
\item
  \textbf{结论}：只要计算时间 \(T(F) \approx T(B)\)，这两个算子就能把彼此的通信时间完全吞掉。
\end{itemize}

\textbf{对冲 B：\(F_{D2}\) 与 \(B_{D1}\) 的完全重叠}

\begin{itemize}
\tightlist
\item
  与上述逻辑镜像，处理反向流的前向和正向流的激活梯度。
\end{itemize}

\textbf{步骤 C：填补 \(W\)（Weight Gradient）}

\begin{itemize}
\tightlist
\item
  \textbf{计算}：上述 \(B\) 计算完成后，立即在后台异步算 \(W\)（权重梯度）。
\item
  \textbf{原理}：\(W\) 不影响任何人，它是``计算废料''，哪里有缝隙就塞在哪里。
\end{itemize}

\Needspace{5\baselineskip}
在一个对冲周期（如 \(F_{D1}\) 与 \(B_{D2}\) 对冲）内共有4个通信步骤，都被掩盖在计算时间中：

\Needspace{5\baselineskip}
在 \(F_{D1}\) 计算时：

\begin{itemize}
\item
  \texttt{Send-Backward-D2}：把上一批次算好的 \(n_{j-1}\) 梯度发给右边。
\item
  \texttt{Recv-Backward-D2}：从左边接收本批次 \(n_j\) 的梯度。
\end{itemize}

\Needspace{5\baselineskip}
在 \(B_{D2}\) 计算时：

\begin{itemize}
\item
  \texttt{Send-Forward-D1}：把算好的 \(m_i\) 激活值发给右边。
\item
  \texttt{Recv-Forward-D1}：从左边接收下一批次 \(m_{i+1}\) 的激活值。
\end{itemize}

\Needspace{5\baselineskip}
\hypertarget{ux4e09ux6536ux5c3eux9636ux6bb5}{%
\subsubsection{（三）收尾阶段}\label{ux4e09ux6536ux5c3eux9636ux6bb5}}

当所有微批次的前向都跑完了，GPU 进入纯反向模式。

\begin{enumerate}
\def\labelenumi{\arabic{enumi}.}
\tightlist
\item
  \textbf{处理剩余的 \(B\)}：把最后几个微批次的输入梯度算完并传走。此时已经没有 \(F\) 可以对冲了，但因为 \(PP\) 被分成了两个方向，气泡的绝对时长减半。
\item
  \textbf{全力算 \(W\)}：这是最后的体力活，所有 GPU 并行计算自己负责的层权重更新。
\end{enumerate}

\Needspace{5\baselineskip}
\hypertarget{ux56dbux6c14ux6ce1ux51cfux5c11ux7684ux6838ux5fc3ux539fux56e0}{%
\subsubsection{（四）气泡减少的核心原因}\label{ux56dbux6c14ux6ce1ux51cfux5c11ux7684ux6838ux5fc3ux539fux56e0}}

在 1F1B 中，Rank 0 算出 \(F_1\) 后，要等 \(2(PP - 1)\) 个时间片才能拿到梯度算 \(B_1\)。

\[
T_{\mathrm{wait}} = (PP - 1) \times (F + B)
\]

\Needspace{5\baselineskip}
DualPipe采取了以下方式：

\begin{enumerate}
\def\labelenumi{\arabic{enumi}.}
\tightlist
\item
\Needspace{5\baselineskip}
  \textbf{双向对冲（Bidirectional）}：

  \begin{itemize}
  \tightlist
  \item
    在 1F1B 等待梯度的时间里，DualPipe 让 GPU 去跑``另一个方向''的前向和反向。
  \item
    \textbf{数学逻辑}：将一个 PP 链拆成两个半深度的 PP 链对冲。气泡从 \((PP - 1)\) 变成了 \(\left(\frac{PP}{2} - 1\right)\)。
  \end{itemize}
\item
\Needspace{5\baselineskip}
  \textbf{\(B\) 与 \(W\) 拆分（Zero Bubble 思想）}：

  \begin{itemize}
  \tightlist
  \item
    在传统做法中，反向传播是 \(B + W\)。必须等这两个都算完，梯度才能发给上一个 GPU。
  \item
    DualPipe 算完 \(B\)（只需矩阵乘法的一半工作量）\textbf{立即就把梯度发走}。这样邻居 GPU 就能提前开始干活。
  \item
    \textbf{结果}：关键路径上的等待时间进一步压缩。
  \end{itemize}
\item
\Needspace{5\baselineskip}
  \textbf{完全重叠（Full Overlap）}：

  \begin{itemize}
  \tightlist
  \item
    利用现代 GPU 的双工带宽（NVLink 可以同时全速 Send 和 Recv）。
  \item
    DualPipe 调度算法保证了：\textbf{当 GPU 在做计算时，带宽始终是满的；当带宽在传输时，SM 核心始终在计算。}
  \end{itemize}
\end{enumerate}

\Needspace{5\baselineskip}
\hypertarget{ux4e94ux6298ux53e0ux6d41ux6c34ux7ebfux6280ux672f}{%
\subsubsection{（五）折叠流水线技术}\label{ux4e94ux6298ux53e0ux6d41ux6c34ux7ebfux6280ux672f}}

\Needspace{5\baselineskip}
DualPipe的``代价''是一个GPU需要存两倍的权重。原因：

\begin{quote}
正向流 Direction 1：从 \(Rank_0 \rightarrow Rank_3\)。Rank 0 必须有第 1 层。

反向流 Direction 2：从 \(Rank_3 \rightarrow Rank_0\)。但请注意，反向流也是一个完整的前向-反向过程。如果它从 Rank 3 开始起步，Rank 3 就必须有第 1 层。

这样一来，第 1 层在 Rank 0 和 Rank 3 都有，参数量自然翻倍了。
\Needspace{5\baselineskip}
为了解决这个问题，可以采用``折叠''的方法：
\end{quote}

\Needspace{5\baselineskip}
在折叠架构下，我们不再让 Rank 0 和 Rank 3 都存第 1 层，而是这样分配 Stage：

\begin{longtable}[]{@{}
  >{\raggedright\arraybackslash}p{\dimexpr\linewidth/3-4\tabcolsep/3\relax}
  >{\raggedright\arraybackslash}p{\dimexpr\linewidth/3-4\tabcolsep/3\relax}
  >{\raggedright\arraybackslash}p{\dimexpr\linewidth/3-4\tabcolsep/3\relax}@{}}
\toprule
物理 Rank & 负责的逻辑 Stage（1x 参数，不冗余） & 物理位置属性 \\
\midrule
\endhead
Rank 0 & Stage 0（入口）和 Stage 7（出口） & 流水线两头 \\
Rank 1 & Stage 1 和 Stage 6 & \\
Rank 2 & Stage 2 和 Stage 5 & \\
Rank 3 & Stage 3 和 Stage 4 & 流水线折返点 \\
\bottomrule
\end{longtable}

\Needspace{5\baselineskip}
虽然每个 Stage 只存了一份（1x 参数），但因为每个 GPU 现在都手握一个靠前的 Stage（如 S0）和一个靠后的 Stage（如 S7），对冲依然可以发生：

\begin{enumerate}
\def\labelenumi{\arabic{enumi}.}
\tightlist
\item
  任务组合：当 Rank 0 正在处理微批次 \(m_{10}\) 的 Stage 0（前向计算）时，它手里可能正好拿着微批次 \(m_1\) 绕了一圈回来的 Stage 7（前向或反向计算）。
\item
\Needspace{5\baselineskip}
  算子对冲：

  \begin{itemize}
  \tightlist
  \item
    计算 A：\(F_{S_0}\)（新批次的前向）。
  \item
    计算 B：\(B_{S_7}\)（老批次的反向）。
  \end{itemize}
\item
\Needspace{5\baselineskip}
  通信掩盖：

  \begin{itemize}
  \tightlist
  \item
    Rank 0 在发送 \(F_{S_0}\) 给 Rank 1 的同时，正好在接收 Rank 1 传回来的 \(B_{S_7}\) 梯度。
  \end{itemize}
\end{enumerate}

这就是``折叠''的精髓：它利用了流水线的深度（一个微批次从 S0 到 S7 走一圈很慢）创造了时间差，让同一个 GPU 在同一时间处理不同生命周期的任务。
但实际上很多大模型并不一定会用``折叠''的方法，原因有两方面。

\Needspace{4\baselineskip}
\begin{enumerate}
\def\labelenumi{\arabic{enumi}.}
\tightlist
\item
  MoE模型的特点
\end{enumerate}

\Needspace{5\baselineskip}
在一个超大规模 MoE 模型中，模型参数被分为两部分：

\begin{itemize}
\tightlist
\item
  \textbf{共享参数（Shared Parameters）}：包括注意力机制（Attention）、层归一化（LayerNorm）和一些共享的稠密层。
\item
  \textbf{专家参数（Expert Parameters）}：MoE 层中成百上千个独立的专家（MLP 块）。
\end{itemize}

这是最核心的工程实现细节：\textbf{专家参数是不参与流水线镜像冗余的。}

\Needspace{5\baselineskip}
在分布式训练中，我们通常同时使用 PP（Pipeline Parallelism）和 EP（Expert Parallelism）：

\begin{enumerate}
\def\labelenumi{\arabic{enumi}.}
\tightlist
\item
  \textbf{PP 决定了路径}：它决定了微批次在哪几个 GPU 之间流转。
\item
  \textbf{EP 决定了存放}：它决定了某一个具体的专家 \(E_{100}\) 物理上存在在哪块显存里。
\end{enumerate}

\Needspace{5\baselineskip}
在 DualPipe 的双向流中：

\begin{itemize}
\tightlist
\item
  当一个微批次运行到某一层需要调用专家时，它会发起一个 All-to-All 通信，去寻找那个专家所在的物理 GPU。
\item
  无论这个微批次是属于``正向流''还是``反向流''，它去寻找的都是同一个物理地址上的专家。
\item
  \textbf{结论：}专家参数在整个集群里依然只存了一份。DualPipe 只需要在每个 Rank 上多存一份轻量级的共享层参数和路由逻辑（Router）即可。
\end{itemize}

\Needspace{4\baselineskip}
\begin{enumerate}
\def\labelenumi{\arabic{enumi}.}
\setcounter{enumi}{1}
\tightlist
\item
  激活值对显存占用更大
\end{enumerate}

对于超长上下文（Context Length）训练，显存瓶颈往往不是参数（Weights），而是激活值。

\begin{itemize}
\tightlist
\item
  在训练时，为了反向传播，每一层的前向计算结果都要存下来。
\item
  DualPipe 的显存压力：真正让显存紧张的是它同时在跑两个方向的流，这意味着显存里要同时存放 \(m\) 批次和 \(n\) 批次的激活值。
\end{itemize}

\Needspace{5\baselineskip}
但 DeepSeek 巧妙地化解了这一点：

\begin{itemize}
\tightlist
\item
  通过 MLA（Multi-head Latent Attention）大幅压缩了 KV Cache 和激活值的体积。
\item
  配合 Recomputation（激活值重计算）技术，可以在需要时再计算部分激活值，用计算时间换取显存空间。
\item
  相比之下，多存那一套共享参数所占用的几 GB 显存，简直是``洒洒水''。
\end{itemize}

\Needspace{4\baselineskip}
\begin{enumerate}
\def\labelenumi{\arabic{enumi}.}
\setcounter{enumi}{2}
\tightlist
\item
  训练调度方面
\end{enumerate}

\Needspace{5\baselineskip}
带宽压榨：

1x 参数的折叠流水线在调度上非常死板（必须走 V 字型）。而 2x 参数的 DualPipe 拥有两条完全独立且对称的路径，这使得它能百分之百跑满双向 NVLink 带宽，调度灵活性极高。

\Needspace{5\baselineskip}
气泡率的极致追求：

在折叠模型中，由于 \(B\) 和 \(W\) 的依赖关系更复杂，很难做到像 DualPipe 那样几乎``零气泡''。DeepSeek 的目标是把 MFU（模型算力利用率）推向 70\% 以上，2x 冗余是一笔划算的交易。

\FloatBarrier
\Needspace{5\baselineskip}
\hypertarget{ux53c2ux8003ux6587ux732e-37}{%
\subsection{参考文献}\label{ux53c2ux8003ux6587ux732e-37}}

\vspace{0.25em}
\begin{itemize}
\tightlist
\item
  Narayanan, D., Harlap, A., Phanishayee, A., et al.~(2019). \href{https://dl.acm.org/doi/10.1145/3341301.3359646}{PipeDream: Generalized Pipeline Parallelism for DNN Training}. SOSP 2019.
\item
  DeepSeek-AI. (2024). \href{https://arxiv.org/abs/2412.19437}{DeepSeek-V3 Technical Report}. arXiv:2412.19437.
\end{itemize}

\clearpage
